\chapter{Сравнение}

\section{Введение}
\section{Сравнение методов профилирования}\label{sec:sravenie}
Для проектирования вентилятора с меридиональным ускорением потока было использовано три подхода:
a) по результатам продувок плоских решёток профилей, игнорируя сужение проточной части.
b) по результатам продувок неподвижных решёток профилей в изолированных тонких осесимметричных слоях течения, образованных линиями тока в меридиональном сечении проточной части. 
c) так же как b, но с учётом вращения решёток профилей.
Вариант a классический. Преимущество подхода b заключено в возможности пересчёта результатов экспериментальных продувок плоских решёток профилей с изменяющейся осевой компонентой скорости с помощью конформных отображений. Переход к неинерционной системе отсчёта при подходе b может привести к ошибкам проектирования ротора. В таком случае подход c окажется предпочтителен несмотря на то, что вряд ли выйдет при этом использовать надёжные экспериментальные данные. 
На одно задание было с использованием перечисленных подходов спроектировано три рабочих колеса вентиляторов. 

Теоретический напор	0,525
Коэффициент производительности		0,6
Сужение площади проточной части		1,14
Относительный диаметр втулки на выходе из проточной части		0,6
Закон закрутки в осевом зазоре	\(C_u R=\const\)

Углы атаки изменялись от минус 2° на периферии до 4 на втулке. Густоты решёток всех роторов одинаковы и определённому по Б—С для ступени с цилиндрической втулкой. Диаметр периферии был постоянен и равен 702 мм. Радиальный зазор составлял 1 мм.
Проектирование осуществлялось с помощью МДВ. Влияние вязкости учитывалось параметром потери циркуляции \(k=\Gamma_\text{вязк}/ \Gamma_\text{невязк} \)
При численном анализе работы ротора на расчётном режиме с невязким рабочим телом (\(k=1\)) получено распределения циркуляции по высоте проточной части (рис. а). Подход c содержит меньше допущений и приводит к результату наиболее близкому к расчётному. Результаты b хорошо соответствуют расчётным у периферии колеса, а по мере приближения к втулке дают завышенное предсказание циркуляции, так как меридиональные линии тока имеют всё больший уклон. Результаты подхода a привели к завышенному предсказанию циркуляции по всей высоте проточной части.
В реальных газах влияние вязкости значительно. Обычно рабочему режиму ротора осевого вентилятора соответствуют значения \(k\) от 1,05 до 1,15. Эти значения можно заранее уточнить для каждого сечения моделированием обтекания решётки профилей вязким рабочим телом. 
В данном случае серией моделирований тонких слоёв течения с моделью турбулентностью k-ω были уточнены значения k на периферийном, среднем и втулочном сечениях. Перепрофилированные рабочие колёса исследовались численным моделированием с использование модели турбулентности \(k-\epsilon\). \(y^{+}=75\). Распределение циркуляции по высоте проточной части на рис. б.  Результат аналогичен случаю с невязким газом.
Опираясь на результаты моделирования предварительного проекта ротора полученным любым из указанных подходов можно уточнить геометрию. Для этого достаточно поправить распределение коэффициента k по высоте проточной части. При этом теряется физический смысл коэффициента, но полученное распределение циркуляции лучше соответствует проектному. В таком случае первоначальным k можно задаваться произвольным в пределах его обычных значений и позже уточнить из результатов моделирования.
На рис. а представлены распределения циркуляции, полученные численным моделированием роторов с заданным k постоянным по всей высоте и равном 1.1. На рис. б циркуляция после перепрофилирования роторов с уточнёнными k для каждого сечения. 
На рис. диаграмма со значениями относительного среднего квадратичного отклонения циркуляции от заданного в пределах относительной высоты проточной части от 0,2 до 0,8 для трёх методов проектирования. Первый столбец означает отклонение при проектировании по заранее определённым значениям k. Второй — при k равном 1,1 по всей высоте проточной части. Третий — при уточнённом k.

Таким образом, все пути проектирования привели к заданным параметрам лопаточной машины в пределах среднего относительного отклонения 3\%. Однако подход c, при котором учитывается изменение кольцевой осесимметричной струйки тока и вращение колеса, приводит к более точным результатам начиная с первой итерации профилирований колёс. Это происходит и в случае заранее уточнённых коэффициентах k для изолированных сечений и для всей высоты проточной части.

\section{Сравнение с вентилятором с цилиндрической проточной частью}

\section{Расширение зоны проектирования}
Для ступени РК+СА с определёнными параметрами \(\bar{H}_\text{т}\)  и \(c_{3a}\) существует степень меридионального ускорения которой достигается максимальный КПД ступени. 
Чем больше \(\bar{H}_\text{т}\) тем выше эта степень. 

Для вентиляторов с цилиндрической проточной частью типичные значения \(\bar{H}_\text{т}\)  от 0,1 до 0,35, а \(c_{3a}\) aот 0,4 до 0,7. Повышение степени меридионального потока может раздвинуть эти границы.
Определим границы, при которых вентиляторы РК+СА имеет КПД равный 87%. 

Границы получены для значений доли динамического напора в теоретическом H"d" /\(\bar{H}_\text{т}\)  от 0,3 до 1. 
При равном КПД доля статического давления связывает напор и производительность, то есть определяет кривую режимов работы вентилятора. 
Вентиляторы, работающие в одном режиме, развивают одинаковое давление и производительность при различных частотах вращения рабочего колеса, чем выше \(\bar{H}_\text{т}\)  и \(c_{3a}\) aтем ниже частота вращения.

Все вентиляторы имели диаметр проточной части 702 мм, радиальный зазор колеса 1 мм. Частота вращения составляла 1000 об/мин. Относительный диаметр втулки на выходе \(d_3\) из проточной части составлял 0,5 0,6 и 0,7. 

Густота решётки профилей в среднем сечении принималась равной густоте цилиндрической ступени, при тех же расчётных параметрах, определённой по Б—С. 
Значение густоты принималось не меньше 0,67 и не больше 2,5.  Длина средней линии профиля оставалась неизменной по высоте лопатки. Углы атаки распределялись линейно по относительной высоте лопатки от минус 2 на периферии до 4 на втулке. 

Ротор и статор моделировались отдельно. КПД ступени определялся как кпдРК+кпдСА-1=КПД. 
Коэффициент потери циркуляции принимался постоянным по высоте лопатки и уточнялся по среднерасходному определённому из предварительномого моделирования течения. При моделировании ротора контролировался полученный \(\bar{H}_\text{т}\) , а при моделировании статора параметр остаточной закрутки \(n_2\) — отклонения от расчётного не превышало 1%.
На каждой кривой режима определялись два вентилятора, c \(\bar{H}_\text{т}\)  второго на 10\% большим первого, такие, чтобы КПД первого был больше 87\%, а второго меньше. 
Значение \(\bar{H}_\text{т}\)  определялось линейной аппроксимацией. 

Принятые выше ограничения существенно повлияли на полученную картину. Особенно принимаемое распределение густоты решётки по высоте проточной части.  В частности, для цилиндрической проточной части не было достигнуто заданное значение КПД на режиме 1, хотя такие схемы существуют с густотой на среднем радиусе меньшей 0,67. Однако принятые ограничения позволили в существенно автоматизировать процесс поиска, и в большинстве случаев приведут к принебрежимым ошибкам.

При относительном диаметре втулки равном 0,5 меридиональное ускорение практически не сдвинуло границы для режимов с долей статики от 0,3 до 0,5. 
Но заметно сдвинуться в области больших напоров можно на режимах статики больше 0,6 при m больше 1. 
При этом \(m_f\) 1,15 позволяет добиться высших напоров, чем при \(m_f\) равном 1,3. 
Однако при режимах, где практически вся подведённая работа преобразуется в скоростной напор, при \(m_f\) 1,3 удаётся сохранить значение КПД 87\% при \(\bar{H}_\text{т}\)  равном 0,29 в противовес 0,24 при \(m_f\) равном 1,15.

В случаях с диаметрами втулок 0,6 и 0,7 наблюдается несколько больший разрыв достигнутых \(Ht\)  при большой доле статического давления, чем для относительно маленького значения относительного диаметра втулки 0,5. При увеличении доли динамического напора заметно разделяются границы для исследованных mf, достигаются максимальные \(\bar{H}_\text{т}\)  при H"d" /\(\bar{H}_\text{т}\)  около значения 0,6. При втулке 0,7 значения \(\bar{H}_\text{т}\) достигают 0,65 при \(m_f\) 1,3. 

\section{Заключение}

\FloatBarrier